\section{La modalidad: Proyecto terminal}
Presentas y defiendes tu proyecto terminal ante tu sínodo. El recorrido es el mismo que el de tesis; los oficios lo nombran «proyecto terminal».

\subsection{Conocer la modalidad y comenzar}
En \emph{Posgrado: obtener mi grado} está la lista de modalidades; cada una abre su página con los requisitos en resumen y los manuales, sin iniciar nada. Para comenzar, acepta el aviso de privacidad y pulsa \boton{Comenzar mi trámite de grado}. La modalidad se elige después, en tus datos.
\captura{proyecto_terminal_egresado_01_entrada}{«Posgrado: obtener mi grado»: los cuatro momentos, las modalidades y sus manuales.}
\captura{proyecto_terminal_egresado_02_informativa}{Página de la modalidad «Proyecto terminal», antes de iniciar: en qué consiste, requisitos en resumen y manuales.}
\captura{proyecto_terminal_egresado_03_comenzar}{Aceptas el aviso de privacidad y pulsas «Comenzar mi trámite de grado».}
\captura{proyecto_terminal_egresado_05_datos}{Tus datos: aquí eliges la modalidad («Proyecto terminal») y capturas los datos de tu director con su correo.}
\section{Lo que cambia en «Proyecto terminal»}
El recorrido es el mismo en todas las modalidades. En los oficios y formatos el trabajo se nombra «proyecto terminal», y en el formato de integración de expediente se marca la casilla de esta modalidad.
\captura{proyecto_terminal_director_01_voto}{Voto aprobatorio de la modalidad «Proyecto terminal»: el documento nombra el trabajo según la modalidad.}
